\BourbakiExerciseGroup

\begin{enumerate}
\def\labelenumi{\arabic{enumi})}
\item
  Let \((x_{i})_{i\in I}\) be a summable family in a complete Hausdorff commutative group \(G\) , and let \(\mathfrak{S}\) be a set of subsets of \(I\) , directed with respect to the relation \(\subset\) and forming a covering of \(I\) . Show that \(\sum_{i\in I}x_{i} = \lim_{J}\sum_{i\in J}x_{i}\) , where the limit is taken with respect to the section filter of \(\mathfrak{S}\) .
\item
  Let G be a complete Hausdorff commutative group such that every neighbourhood of o in G contains an open subgroup of G (this is the case in particular if G is locally compact and totally disconnected; cf.~§ 4, Exercise 18). A family \((x_{t})_{t\in I}\) of points of G is summable if and only if \(\lim x_t = 0\) with respect to the filter of complements of finite subsets of I. Hence show that every convergent series in G is commutatively convergent.
\item
  Let \((x_{i})_{i\in I}\) be a family of points of a Hausdorff commutative group G. For every finite subset H of I, let \(s_{H}\) denote \(\sum_{i\in H}x_{i}\) , and
\end{enumerate}

let A be the set of cluster points of the mapping \(H \rightarrow s_{H}\) with respect to the directed set \(\mathfrak{F}(I)\) . For every finite subset J of I, let \(\Phi(J)\) denote the closure of the set of points \(s_{H}\) , where H runs through the finite subsets of I such that \(H \cap J = \varnothing\) . Show that if \(A \neq \varnothing\) then the set A - A (i.e.~the set of all \(x - x'\) where \(x \in A\) and \(x' \in A\) ) is the same as the set \(B = \bigcap_{J \in \mathfrak{F}(I)} \Phi(J)\) ; and that in any case \(B + B \subset B\) . Deduce that if A is not empty, then B is a closed subgroup of G and that A is a coset of B in G.

\begin{enumerate}
\def\labelenumi{\arabic{enumi})}
\setcounter{enumi}{3}
\item
  With the notation of Exercise 3, take G to be the discrete additive group Z of rational integers and I to be the set N. Give an example of a non-summable sequence \((x_{n})\) for which the set A consists of the single point o (choose the \(x_{n}\) so that every finite partial sum, all of whose terms have indices \(\geqslant m\) , is an integral multiple of m).
\item
  Let \((x_{n})\) be a sequence of points of a Hausdorff commutative group. If, for every infinite subset I of N, the series defined by the sequence \((x_{n})_{n\in I}\) is convergent, then the sequence \((x_{n})\) is summable (argue by reductio ad absurdum as in the proof of Proposition 9).
\end{enumerate}

¶ 6) a) Let \(\sigma\) be a permutation of \(\mathbf{N}\) and let \(\varphi(n)\) be the smallest number of intervals of \(\mathbf{N}\) whose union is \(\sigma([0, n])\) . Suppose that \(\varphi(n)\) is bounded as \(n\) runs through \(\mathbf{N}\) . Then if \((u_n)\) is any convergent series whose terms belong to a Hausdorff commutative group \(\mathbf{G}\) , the series \((u_{\sigma(n)})\) is convergent and has the same sum as \((u_n)\) .

* b) Suppose that \(\varphi(n)\) is not bounded in N. Construct a series \((u_n)\) whose terms belong to the additive group R and which converges in R but is such that the series \((u_{\sigma(n)})\) does not converge. {[}Consider a strictly increasing sequence \((m_k)\) of integers, defined by induction on \(k\) , satisfying the following conditions: (i) if \([0, n_k]\) is the largest interval of N with left-hand extremity o which is contained in \(\sigma([0, m_k])\) , then

\[
\sigma ([ 0, m _ {k} ]) \subset [ 0, n _ {k + 1} ];
\]

\begin{enumerate}
\def\labelenumi{(\roman{enumi})}
\setcounter{enumi}{1}
\tightlist
\item
  \(\varphi(m_k) \geqslant k + 1\) . Then define \(u_n\) appropriately for \(n_k < n < n_{k+1}]\) . Generalize to the case of a complete Hausdorff commutative group G which is generated by every neighbourhood of o. *
\end{enumerate}

\begin{enumerate}
\def\labelenumi{\arabic{enumi})}
\setcounter{enumi}{6}
\tightlist
\item
  Let \((x_{mn})\) be a double sequence of points of a Hausdorff commutative group \(\mathbf{G}\) , satisfying the following conditions:
\end{enumerate}

\begin{enumerate}
\def\labelenumi{(\alph{enumi})}
\item
  The series defined by the sequence \((x_{mn})_{n\in \mathbb{N}}\) is convergent for each \(m\geqslant 0\) ; let \(y_{m}\) be its sum.
\item
  If we put \(r_{mn} = \sum_{p=n}^{\infty} x_{mp}\) , then the series defined by the sequence \((r_{mn})_{m \in \mathbb{N}}\) is convergent for each \(n \geqslant 0\) ; let \(t_n\) be its sum.
\end{enumerate}

Show that for each \(n \geqslant 0\) the series defined by the sequence \((x_{mn})_{m \in \mathbb{N}}\) is convergent, and let \(z_{n}\) be its sum. The series whose general term is \(y_{m}\) has the same sum as the series whose general term is \(z_{n}\) if and only if \(t_{n}\) tends to 0 as n tends to infinity.

